\section[title={2024-09-16 - Figurenkonstillation},reference={2024-09-16-Figurenkonstillation}]

\startsectionlevel[title={Figuren},reference={figuren}]

\startsectionlevel[title={Vater, ein
\quotation{Tyran}},reference={vater-ein-tyran}]

\margintext{{\bf Egozentrisch:} die eigene Person als Zentrum allen Geschehens
betrachtend; alles in Bezug auf die eigene Person beurteilend und eine
entsprechende Haltung erkennen lassend.
}

\startitemize[packed]
\item
  egozentrisch
  \startitemize[packed]
  \item
    Eigenes Welt- / Rollenbild
  \item
    Zwingt dieses anderen auf
    \startitemize[packed]
    \item
      Beispiel: Nötigt die Familie zusammen Musik zu hören
    \stopitemize
  \stopitemize
\item
  Konsequenz bei Wiederspruch
  \startitemize[packed]
  \item
    bis zur gewaltbereitschaft
  \stopitemize
\item
  Weist stets Schuld von sich sein Familienbild ist ausschlaggebend;
  zeigt kaum {\bf Emotionen}
\stopitemize

\stopsectionlevel

\startsectionlevel[title={Mutter, die
\quotation{Unsichere}},reference={mutter-die-unsichere}]

\startitemize[packed]
\item
  Wird von Vater unterdrückt
\item
  hasst Streit, weswegen sie sich auch nicht trennt
\item
  Lehrerin, eigene Kinder nehmen sie nicht als böse oder streng war. Die
  Schüler aber schon
\item
  Gefühlsmensch; spielt heimlich im Schrank Geige
\item
  \quotation{verpätzt} die Kinder an den Vater
\stopitemize

\stopsectionlevel

\startsectionlevel[title={Tochter, die
\quotation{Rebelling}},reference={tochter-die-rebelling}]

\startframedtask
\startitemize[packed]
\item
  Seiten: 35, 24, 46, 99, 30
\item
  S. 30
\stopitemize

\stopframedtask

\startitemize[packed]
\item
  hält Stille aus
\item
  hat kein Spaß am Muschelnessen
\item
  reflektiert / hintergragt Familienstrukut
\item
  belügt Familie, um Freiraum zu haben
\item
  bezeichnet Mutter als Petze
\item
  Verurteilt Verhalten der Mutter, bei anwesenheit des Vaters
\stopitemize

\stopsectionlevel

\startsectionlevel[title={Sohn, der
\quotation{Sensible}},reference={sohn-der-sensible}]

\startitemize[packed]
\item
  Will Vater stolz machen
\item
  Wird vom Vater zur \quotation{Härte} erziehen
\item
  Suizidal
\item
  passt sich dem Vater an
\item
  ähnelt Mutter in Emotionalität
\stopitemize

\stopsectionlevel

\stopsectionlevel

\startsectionlevel[title={Eine \quotation{richtige}
Familie},reference={eine-richtige-familie}]

\startframedtask
{\bf Stellen} Sie die Vorstellung des Vaters von einer
\quotation{richtigen} Familie mithilfe der Seiten: 10f., 21, 38f., 49,
62f., 74f., 77ff., 93, 98f., 103ff. und 107f. {\bf dar}.

\stopframedtask

\startitemize[packed]
\item
  Familie verhält sich während anwesenheit des Vaters anders
\item
  Strenges Regelsystem durch Vater
\item
  Vater gestalltet Leben der Familie um sein eigenes
\item
  Vater ist Chef der Famiiie (verkörpert traditionelle Werke)
\item
  Vater ist arm aufgewachsen, will Armut vermeiden
\item
  Traditionelle Rollenverteilung. Er bringt das Geld, die Mutter kömmert
  sich um den Haushalt
\item
  Zusammenhalt der Ehe im Voredergrund
\stopitemize

\stopsectionlevel
